\section*{写在前面：这份讲义是什么，不是什么}

先花三分钟，把话说清楚。

这份讲义是写给谁的？写给成人 ADHD、自认 ADHD、以及正怀疑自己可能是 ADHD 的成年人的。你不需要诊断书，也不需要用「没确诊」来否定自己——这本讲义讨论的是「这套大脑怎么运转」，诊断与否，不影响它对你有多少用。

这份讲义从哪来？它的思想内核，提炼自 B 站 UP 主「阿卡迪萨」（ADHD 野路子）的 24 期视频。一个清华数学博士，在 B 站搞笑的过程中发现自己疑似 ADHD，然后把自己的大脑当成一台机器来拆、来猜、来调试，留下了一整套自成一派的猜想。这份讲义的最大特点，正是继承了那个态度：\term{不讲医学权威，只讲自己大脑的实际运行日志}。

那这份讲义不是什么？它不是医学诊断，也不是吃药指南。它不承诺治愈。它只做一件事：\term{让你第一次觉得，自己没坏}。

\key{最舒适的认识，一句话：你不需要把自己修成正常人。你需要的是读懂自己的参数，然后顺着参数生活。}

所谓「舒适」，不是躺平。舒适是停止和重力对抗——你不再天天质问自己「为什么别人能做到我却不能」，而是开始研究「我这套系统到底吃什么燃料、走什么路线、在什么条件下运转得最好」，然后照着这个去设计生活。

讲义的路线图，一共四部分：

\begin{checklist}
  \item \textbf{底层认知}：ADHD 到底是一套什么样的算法，它真正的麻烦在哪；
  \item \textbf{理论延伸}：同一个算法，为什么在学历上层和学历中下层会走出两种人生；
  \item \textbf{自救方法论}：先让认知舒服，再给一套可操作的方法库，最后从你的位置出发；
  \item \textbf{更深一层的思考}：一个更舒服的看待世界、看待自己的角度，和一页纸的每日清单。
\end{checklist}

文末还有两份附录：一份给「怀疑自己是 ADHD」的人的自检线索，一份本讲义的黑话速查表，方便你随时回来翻。

最后，一个必要的声明，也是原视频反复强调的态度：\term{以下是猜测，个体差异极大}。所有方法都是「试试看」，不是「必须做」。好用就用，不好用就扔。这份讲义里的每一个办法，原样照搬大概率会失灵，按自己的参数改装之后才可能灵。

如果你已经被生活严重困住，该寻求专业帮助就去找，这不丢人，也和读这份讲义不冲突。但在这之前、在这之后，都请先记住这句话——你不是懒，你只是还没被读懂。
